%!TEX root = ../main.tex
% =====================================================================
% Related work for "Know-When-to-Trust Crop-Row Guidance".
% Honest framing: the contribution is a TRUST LAYER (conformal heading
% intervals + risk-controlled steering policy) on a STOCK detector with a
% robust Huber-IRLS line fit. We organise prior art into five threads and
% end each with the explicit gap this paper fills:
%   (a) vision crop-row guidance (detect-then-fit, equal/heuristic weights);
%   (b) lane / row-anchor detection (the row-line paradigm we borrow);
%   (c) conformal + selective prediction, INCLUDING conformal-on-detector
%       prior art we CONTRAST so we never claim "first conformal";
%   (d) detection box-uncertainty (DFL/GFL) -- the NEGATIVE control here;
%   (e) edge / INT8 deployment.
% No "first/novel" hype. The DFL-variance precision-weighting thesis is the
% dead negative control, never a contribution.
% =====================================================================

\section{Related work}
\label{sec:related}

\noindent We organise prior art around five threads that bound our
contribution. The first two cover the application and the line-detection
paradigm we reuse: vision-based crop-row guidance and the lane/row-anchor
formulation it descends from. The third is the statistical machinery our
trust layer instantiates---conformal and selective prediction---including the
recent conformal-on-detector work we explicitly contrast against, so that we
make no claim of being first to apply conformal prediction to a detector. The
fourth is detection box-uncertainty (DFL/GFL), which we examine and report as a
\emph{non-informative} signal for heading error here. The fifth is edge and
integer-quantised deployment, the setting in which the guidance must run. After
each thread we state explicitly what it leaves open, so that the gap this paper
fills is unambiguous.

\subsection{Vision-based crop-row guidance}
\label{subsec:rw_nav}
\noindent Guiding a robot or tractor down a crop row from a forward-facing
camera is a long-standing task; early vision-guided tractor systems date back
decades~\cite{cth6}, and classical pipelines recover row geometry without
learned segmentation, e.g.\ intensity-based crop-row determination by image
analysis~\cite{sogaard} and global energy minimisation over candidate row
models~\cite{crbd}. Modern systems converge on a small number of output
paradigms grouped by how the guidance line is produced.

\paragraph{Detect-then-fit.}
The de-facto recipe localises row segments or plant clusters with an object
detector and fits a single guidance line through the box centroids. The
YOLOX-based maize-row system of Gong et al.~\cite{diao_yolox} is
representative: it detects row patches, takes their centroids, and fits a
centreline by least squares, reporting angular error. This detect-then-fit
recipe with equal or heuristic point weights is the lineage of our equal-weight
least-squares baseline (E1), and it is the family our stock detector belongs
to. Real-time computer-vision crop-row detection on robots follows the same
template~\cite{cth3}.

\paragraph{Segmentation plus geometric scan.}
A second paradigm segments the crop/soil mask and recovers the row geometry by a
scan-line or triangle search over the mask. The de Silva et al.\ line of
work~\cite{cth1,cth2} validates this across many field conditions and,
importantly, reports results in navigation units (cross-track and heading
error) on the CRDLD benchmark~\cite{cth1}; we adopt its mask-derived centreline
as the detector-independent ground-truth line for in-domain evaluation. The
geometric scan itself is a hard, non-differentiable search and carries no
per-frame reliability estimate.

\paragraph{End-to-end row regression.}
A third paradigm regresses the row directly. CCRDNet extracts the central crop
row for universal in-row navigation~\cite{cth11}; E2CropDet offers an efficient
end-to-end solution~\cite{cth12}; RowDetr regresses the row as a per-row
polynomial with INT8 edge deployment~\cite{rowdetr}; and row--column attention
recasts the pipeline as an end-to-end attention model~\cite{rowcolattn}. These
methods produce the curve parameters directly, but with no per-frame uncertainty
on the recovered line.

\paragraph{What this thread leaves open.}
Two observations recur and motivate our work. First, the literature increasingly
evaluates in \emph{navigation units}---heading error in degrees and cross-track
error---rather than detection mAP, because a detector that scores well on mAP
can still yield a poor guidance line~\cite{cth1,cth2,diao_yolox}; a method that
aims to improve navigation must be measured there. Second, and central to us:
across all paradigms the guidance line is emitted as a \emph{point} estimate.
The line is fit with equal or heuristic weights, and the system reports a single
heading---at best gated by a hand-tuned detection-confidence threshold---with no
\emph{coverage guarantee} and no principled signal for \emph{when not to trust
it}. Vision-only field navigation without RTK-GPS needs exactly that: a guidance
line \emph{and} a calibrated decision of when to abstain. This is the gap our
trust layer (C2) fills, on top of a robust point readout (C1) that uses Huber-IRLS rather than equal-weight least squares.

\subsection{Lane and row-anchor detection}
\label{subsec:rw_lane}
\noindent The row-line formulation used in agriculture descends from lane
detection. Ultra-Fast structure-aware lane detection casts the problem as
row-anchor classification---selecting a horizontal location per image
row~\cite{ufld}---and LaneATT performs anchor-based real-time lane
detection~\cite{laneatt}. These ideas have been ported to agricultural rows,
e.g.\ the anchor-line network ALNet~\cite{alnet} and row--column
attention~\cite{rowcolattn}. Closest to our point-readout stage, Van Gansbeke et
al.\ fit a lane as a differentiable least-squares line over points emitted by a
detector~\cite{vangansbeke}---the detect-then-fit-a-line precedent we build on.

\paragraph{What this thread leaves open.}
These methods are accurate but the per-anchor or per-point confidence, where
present, is a classification score, not a calibrated statement about the error
of the recovered line. The differentiable least-squares fit of~\cite{vangansbeke}
weights points but provides no finite-sample interval on the resulting heading.
None of these methods emit a heading prediction \emph{interval} with a target
coverage, nor a reject option tied to that interval---which is what downstream
steering on low-cost hardware requires.

\subsection{Conformal prediction, selective prediction, and conformal on detectors}
\label{subsec:rw_conformal}
\noindent Our trust layer is an instantiation of conformal prediction. Conformal
methods produce prediction sets with finite-sample, distribution-free coverage
under exchangeability~\cite{vovk2005}; split (inductive) conformal prediction
makes this practical for regression by calibrating a nonconformity score on a
held-out set and reading off a quantile to form intervals~\cite{lei2018}, as
summarised in recent tutorials~\cite{angelopoulos2023}. Conformalized quantile
regression yields \emph{locally adaptive} interval widths by scaling the
conformal quantile with a per-input scale~\cite{romano2019cqr}; our score-scaled
half-width $h=\hat{q}\,s$ (Section~\ref{sec:related} notation; the
nonconformity score $s$ is a leave-one-out heading dispersion) is in this
locally adaptive family. The accompanying abstain/slow/proceed policy is a
selective-prediction mechanism: prediction with a reject
option~\cite{geifman2017selective}, made distribution-free by conformal risk
control~\cite{angelopoulos2022rcps}, which is the formal basis of our
risk-controlled gate (E4, E6).

\paragraph{Conformal applied to object detectors (prior art we contrast).}
Conformal prediction has recently been applied directly to object-detection
outputs, producing coverage-controlled boxes or detection
uncertainty~\cite{confdet1,confdet2}. We therefore do \emph{not} claim primacy
in pairing conformal prediction with a detector. The distinction we draw is in
\emph{what} is made conformal and \emph{for what task}: prior
conformal-on-detector work calibrates the detector's own outputs (e.g.\
box coordinates or class/objectness) for the detection task, whereas we calibrate
a \emph{navigation-native, geometric} quantity---the heading of a robustly fit
guidance line---using a label-free, failure-mode-matched score (leave-one-out
heading dispersion) that targets wrong-row contamination, the dominant guidance
failure. We position our method as a \emph{domain instantiation} of these ideas
for crop-row steering, not as a new conformal method.

\paragraph{What this thread leaves open.}
The coverage guarantee of split conformal is finite-sample-valid only under
exchangeability \emph{within one distribution}~\cite{vovk2005,lei2018}. Crop-row
guidance must run across fields, crops, and acquisition conditions, where
exchangeability is violated; we therefore treat cross-dataset (CRDLD$\leftrightarrow$CRBD)
and field-video transfer as a coverage-\emph{degradation} study (E2), not as a
guarantee. Bringing a label-free conformal heading interval and a
risk-controlled steering policy to a stock edge detector---and quantifying how
its coverage degrades off-distribution and after integer quantisation---is the
contribution this thread sets up.

\subsection{Box-distribution uncertainty in detectors (the negative control)}
\label{subsec:rw_unc}
\noindent Modern stock YOLO detectors~\cite{ct4,yolov10} regress each box edge
with a Distribution Focal Loss (DFL) head from Generalized Focal
Loss~\cite{gfl}: instead of a single value, each edge $e$ predicts a softmax
distribution $p_e(k)$ over bins, whose mean is the edge offset and whose
variance $\sigma^2_{\mathrm{DFL}}=\sum_k (k-\mu_e)^2\,p_e(k)$ encodes the
network's spread. GFLv2 reads the \emph{shape} of this distribution into a
scalar localisation-quality score used \emph{inside} the detector for ranking
and NMS~\cite{gflv2}. A parallel line attaches explicit localisation
uncertainty: KL-Loss predicts per-coordinate variance for variance-voting
NMS~\cite{klloss}, Gaussian-YOLOv3 attaches a localisation-uncertainty output to
each box~\cite{gaussianyolo}, and Bayesian deep learning separates aleatoric and
epistemic uncertainty more generally~\cite{kendallgal}.

\paragraph{What this thread leaves open---and why it fails here.}
It is tempting to read $\sigma^2_{\mathrm{DFL}}$, already produced for free by
the stock head, as a per-detection reliability and to precision-weight the line
fit by its inverse. We tested exactly this hypothesis and report it as a
negative control (E0): the DFL box-distribution variance is \emph{not}
informative of realised heading error in our setting (Spearman
$\rho=-0.13$). The dominant guidance failure is not a slightly blurry box but a
detection on the \emph{wrong} row, which the per-box DFL spread does not flag.
Consequently we do \emph{not} use $\sigma^2_{\mathrm{DFL}}$ (nor a KL-Loss or
Gaussian-YOLO branch~\cite{klloss,gaussianyolo}) as the weight or the trust
signal. Robustness, not variance weighting, is what works for the point fit
(Huber-IRLS; C1, E1), and a failure-mode-matched leave-one-out dispersion, not
$\sigma^2_{\mathrm{DFL}}$, is what carries the trust signal (C2, E0/E3). This
negative result is what motivates the rest of the paper.

\subsection{Edge and integer-quantised deployment}
\label{subsec:rw_edge}
\noindent Vision-only guidance for low-cost robots must run on embedded
accelerators. The efficiency lineage includes compact backbones
(MobileNet/MobileNetV2~\cite{ct16,ct19}, EfficientNet~\cite{ct20}), knowledge
distillation~\cite{ct18}, and structural pruning (pruning filters for efficient
convnets~\cite{ct17}, DepGraph~\cite{cth14}). Deployment on the target hardware
relies on integer-arithmetic inference: post-training INT8 quantisation in the
sense of Jacob et al.~\cite{jacob2018int8}, NVIDIA TensorRT for Jetson-class
FP16/INT8 inference~\cite{tensorrt}, and the Rockchip RKNN toolkit for INT8 NPU
deployment~\cite{rknntoolkit}.

\paragraph{What this thread leaves open.}
This literature optimises detector latency, footprint, and detection accuracy
under quantisation, but does not ask what happens to a \emph{statistical coverage
guarantee} when the detector is quantised to INT8. Because our trust layer is
CPU-side and post-NMS, it adds no operator to the feature maps or NPU---we state
this plainly, not as a selling point: the heading readout and conformal interval
run after detection on the CPU. The open question we answer (E5) is empirical:
does the realised coverage of the heading interval survive INT8 quantisation on
Jetson Nano (TensorRT) and on an RKNN NPU, and at what latency and throughput?

\subsection{Positioning}
\label{subsec:rw_position}
\noindent In summary, this paper is an \emph{applied} contribution, not a methods
breakthrough and not a new detector. On top of a stock detector~\cite{ct4} with a
robust Huber-IRLS line fit (C1)---which in our experiments matches a RANSAC baseline and improves on
equal-weight least squares, and on which the native DFL
variance~\cite{gfl} is shown to be non-informative (E0)---we add a label-free,
failure-mode-matched conformal heading interval and a risk-controlled steering
policy (C2), positioned as a domain instantiation of conformal and selective
prediction~\cite{vovk2005,lei2018,romano2019cqr,geifman2017selective,angelopoulos2022rcps}
and contrasted against, rather than claiming primacy over, prior
conformal-on-detector work~\cite{confdet1,confdet2}. We evaluate it in navigation
units, across datasets (CRDLD$\leftrightarrow$CRBD~\cite{cth1,crbd}) and on real
maize field video, and on edge hardware with INT8
quantisation~\cite{tensorrt,rknntoolkit,jacob2018int8} (C3). Coverage claims are
finite-sample-valid only under exchangeability within one distribution; off-distribution
behaviour is reported as a degradation study, and the policy replay (E6) is
open-loop over a kinematic bicycle model~\cite{rajamani2012}, not a closed-loop
field run.
